\chapter{अक्षीय सममिति}\label{ch:g6:symmetry}

काग़ज़ की शीट को किसी रेखा पर मोड़ो: हर बिंदु दूसरे बिंदु पर आ
बैठता है, उसका दर्पण-प्रतिबिंब। यही मोड़ना \emph{\omterm{def:g6:symmetry:def}{अक्षीय सममिति}}
(या परावर्तन) है, इस पुस्तक का पहला ज्यामितीय रूपांतरण --- दूसरा,
केंद्रीय सममिति, \cref{ch:g7:central} में आता है।

\section{परावर्तन}

\begin{definition}[सममित बिंदु]\label{def:g6:symmetry:def}
मान लो $d$ एक रेखा है। $d$ के पार किसी बिंदु $M$ का
\emph{सममित}\index{अक्षीय सममिति} बिंदु वह $M'$ है जिसके लिए $d$,
$[MM']$ का \omterm{def:g6:symmetry:bisector}{लंब समद्विभाजक} हो: \omterm{def:g6:lines:objects}{रेखाखंड} $[MM']$, $d$ पर \omterm{def:g6:lines:perp}{लंब} है और
$d$ उसे उसके मध्यबिंदु पर काटती है। $d$ \emph{पर} पड़ा बिंदु
अपना ही सममित होता है।
\end{definition}

\begin{omfigure}
\begin{tikzpicture}[scale=1.05]
  \draw[thick, omExo] (0,-1.7) -- (0,1.9) node[above, font=\small] {$d$};
  \fill (-2.1,1.0) circle (1.6pt) node[left, font=\small] {$M$};
  \fill (2.1,1.0) circle (1.6pt) node[right, font=\small] {$M'$};
  \draw[dashed] (-2.1,1.0) -- (2.1,1.0);
  \draw (-0.10,0.90) -- (0.10,1.10);
  \draw (-1.15,0.90) -- (-0.95,1.10);
  \draw (0.95,0.90) -- (1.15,1.10);
  \draw[thin] (0,1.0) ++(0.16,0) -- ++(0,0.16) -- ++(-0.16,0);
  \fill (-1.3,-1.1) circle (1.6pt) node[left, font=\small] {$N$};
  \fill (1.3,-1.1) circle (1.6pt) node[right, font=\small] {$N'$};
  \draw[dashed] (-1.3,-1.1) -- (1.3,-1.1);
  \fill (0,0.2) circle (1.6pt) node[above right, font=\small] {$P = P'$};
\end{tikzpicture}

{\small रेखा $d$ पर परावर्तन: अक्ष $[MM']$ पर \omterm{def:g6:lines:perp}{लंब} है और उसके
मध्यबिंदु से होकर जाती है (बराबर लकीरें)। अक्ष पर पड़ा बिंदु $P$
हिलता नहीं।}
\end{omfigure}

\begin{method}[सममित बिंदु बनाना]\label{met:g6:symmetry:construct}
रेखा $d$ के पार $M$ का \omterm{def:g6:symmetry:def}{सममित} बिंदु $M'$ बनाने के लिए:
\begin{enumerate}
  \item $M$ से होकर $d$ पर \omterm{def:g6:lines:perp}{लंब} खींचो (\omterm{def:g6:shapes:triangles}{समकोण} पट्टिका); जहाँ वह
        $d$ को काटे उसे $H$ कहो;
  \item परकार या मापपट्टी से $MH$ नापो;
  \item $M'$ को उसी \omterm{def:g6:lines:perp}{लंब} पर, $d$ की दूसरी ओर रखो, जहाँ
        $HM' = MH$ हो।
\end{enumerate}
जाल वाले काग़ज़ पर यह तेज़ है: $M$ से अक्ष तक के खाने गिनो, और
दूसरी ओर उतने ही गिन लो।
\end{method}

\begin{example}[जाल पर]\label{ex:g6:symmetry:grid}
अक्ष जाल की कोई खड़ी रेखा हो और $M$ उससे $3$ खाने बाईं ओर, तो
$M'$ उससे $3$ खाने दाईं ओर, उसी पंक्ति में होता है। तिरछी अक्ष के
लिए गिनती विकर्णों पर चलती है --- हमेशा अक्ष पर \omterm{def:g6:lines:perp}{लंब}।
\end{example}

\begin{omfigure}
\begin{tikzpicture}[scale=0.55]
  \draw[gray!40, thin] (0,0) grid (12,6);
  \draw[thick, omExo] (6,-0.4) -- (6,6.4);
  \draw[thick, omDef, fill=omDef!15]
    (1,1) -- (4,1) -- (4,3) -- (3,3) -- (3,5) -- (1,5) -- cycle;
  \draw[thick, omThm, fill=omThm!15]
    (11,1) -- (8,1) -- (8,3) -- (9,3) -- (9,5) -- (11,5) -- cycle;
  \node[omDef, font=\small] at (2.4,0.45) {आकृति};
  \node[omThm, font=\small] at (9.6,0.45) {उसका दर्पण-प्रतिबिंब};
\end{tikzpicture}

{\small खड़ी अक्ष पर एक आकृति और उसका परावर्तन: हर \omterm{def:g5:solids:def}{शीर्ष}
खाना-दर-खाना परावर्तित होता है, और बायाँ-दायाँ आपस में बदल जाते हैं।}
\end{omfigure}

\begin{proposition}[परावर्तन क्या बचाए रखता है]\label{prop:g6:symmetry:preserved}
किसी आकृति का, किसी रेखा के पार, \omterm{def:g6:symmetry:def}{सममित} प्रतिबिंब \emph{उसी
आकार और उसी नाप} की आकृति होता है: लंबाइयाँ, कोण, \omterm{def:g4:measure:perimeter}{परिमाप} और
\omterm{def:g4:measure:area}{क्षेत्रफल} सब बचे रहते हैं। \omterm{def:g6:lines:objects}{रेखाखंड} का \omterm{def:g6:symmetry:def}{सममित} \omterm{def:g6:lines:objects}{रेखाखंड} होता है,
वृत्त का उसी \omterm{def:g3:shapes:shapes}{त्रिज्या} वाला वृत्त, रेखा का रेखा।
\end{proposition}
\admitted

\begin{remark}
एक चीज़ \emph{नहीं} बचती: दिशा-क्रम। परावर्तित अक्षर $b$, अक्षर
$d$ बन जाता है --- दर्पण-लेखन। अपनी आकृति उतारकर अनुरेखण काग़ज़ पलटने पर
वही प्रतिबिंब बने, तो परावर्तन सही है।
\end{remark}

\section{किसी आकृति की सममिति अक्ष}

\begin{definition}[सममिति अक्ष]\label{def:g6:symmetry:axis}
कोई रेखा $d$ किसी आकृति की \emph{सममिति अक्ष}\index{सममिति अक्ष}
तब होती है जब $d$ पर परावर्तन आकृति को ठीक उसी पर भेज दे ---
$d$ पर मोड़ने से दोनों आधे मिल जाएँ।
\end{definition}

\begin{omfigure}
\begin{tikzpicture}[scale=0.62]
  % समद्विबाहु त्रिभुज
  \draw[thick] (0,0) -- (2.4,0) -- (1.2,2.6) -- cycle;
  \draw[omThm, dashed] (1.2,-0.5) -- (1.2,3.0);
  \node[below, font=\small] at (1.2,-0.7) {समद्विबाहु: $1$ अक्ष};
  % आयत
  \begin{scope}[xshift=5.6cm]
  \draw[thick] (0,0.3) rectangle (3.2,2.3);
  \draw[omThm, dashed] (1.6,-0.4) -- (1.6,3.0);
  \draw[omThm, dashed] (-0.5,1.3) -- (3.7,1.3);
  \node[below, font=\small] at (1.6,-0.7) {आयत: $2$ अक्ष};
  \end{scope}
  % वर्ग
  \begin{scope}[xshift=12.0cm]
  \draw[thick] (0.3,0.1) rectangle (2.7,2.5);
  \draw[omThm, dashed] (1.5,-0.5) -- (1.5,3.1);
  \draw[omThm, dashed] (-0.3,1.3) -- (3.3,1.3);
  \draw[omThm, dashed] (0.0,-0.2) -- (3.0,2.8);
  \draw[omThm, dashed] (0.0,2.8) -- (3.0,-0.2);
  \node[below, font=\small] at (1.5,-0.7) {वर्ग: $4$ अक्ष};
  \end{scope}
\end{tikzpicture}

{\small जानी-पहचानी आकृतियों की \omterm{def:g6:symmetry:axis}{सममिति अक्ष}। आयत से सावधान:
उसके विकर्ण \omterm{def:g6:symmetry:axis}{सममिति अक्ष} \emph{नहीं} हैं (किसी विकर्ण पर मोड़ो ---
वर्ग न हो तो आधे नहीं मिलते)।}
\end{omfigure}

\begin{example}\label{ex:g6:symmetry:shapes}
\omterm{def:g6:symmetry:axis}{सममिति अक्ष} गिनना:
\begin{itemize}
  \item \omterm{def:g6:shapes:triangles}{समद्विबाहु त्रिभुज}: $1$ (ऊपर के \omterm{def:g5:solids:def}{शीर्ष} और आधार के बीच से
        होकर); \omterm{def:g6:shapes:triangles}{समबाहु त्रिभुज}: $3$;
  \item आयत: $2$ (दोनों ``बीच वाली रेखाएँ'', विकर्ण \emph{नहीं});
        \omterm{def:g6:shapes:quadrilaterals}{समचतुर्भुज}: $2$ (उसके विकर्ण!); वर्ग: $4$;
  \item वृत्त: केंद्र से होकर जाने वाली हर रेखा --- अनगिनत।
\end{itemize}
\end{example}

\begin{example}[सममिति बराबरी सिद्ध कर देती है]\label{ex:g6:symmetry:proves}
\omterm{def:g6:shapes:triangles}{समद्विबाहु त्रिभुज} के दोनों आधार-कोण बराबर क्यों होते हैं?
त्रिभुज को उसकी \omterm{def:g6:symmetry:axis}{सममिति अक्ष} पर मोड़ो: बायाँ \omterm{def:g3:division:half}{आधा} ठीक दाएँ आधे पर
आ बैठता है, इसलिए बायाँ आधार-कोण दाएँ आधार-कोण पर आता है ---
दोनों की नाप एक ही होनी चाहिए
(\cref{prop:g6:symmetry:preserved}: परावर्तन कोण बचाए रखते हैं)।
\end{example}

\section{लंब समद्विभाजक}

\begin{definition}[लंब समद्विभाजक]\label{def:g6:symmetry:bisector}
किसी \omterm{def:g6:lines:objects}{रेखाखंड} $[AB]$ का \emph{लंब समद्विभाजक}\index{लंब समद्विभाजक}
वह रेखा है जो $[AB]$ पर \omterm{def:g6:lines:perp}{लंब} हो और उसके मध्यबिंदु से होकर जाए ---
यही वह \omterm{def:g6:symmetry:axis}{सममिति अक्ष} है जो $A$ और $B$ की अदला-बदली करती है।
\end{definition}

\begin{proposition}[समदूरी]\label{prop:g6:symmetry:bisector}
कोई बिंदु $M$, $[AB]$ के \omterm{def:g6:symmetry:bisector}{लंब समद्विभाजक} पर ठीक तभी होता है जब वह
$A$ और $B$ से एक ही दूरी पर हो: $MA = MB$।
\end{proposition}
\admitted

\begin{method}[परकार से रचना]\label{met:g6:symmetry:bisector}
बिना नापे $[AB]$ का \omterm{def:g6:symmetry:bisector}{लंब समद्विभाजक} खींचने के लिए:
\begin{enumerate}
  \item परकार को $AB$ के आधे से ज़्यादा खोलो;
  \item $A$ को केंद्र मानकर \omterm{def:g6:lines:objects}{रेखाखंड} के ऊपर और नीचे दो चाप खींचो;
  \item \emph{उसी} खुलाव से $B$ को केंद्र मानकर दो चाप खींचो;
  \item दोनों कटान-बिंदुओं को जोड़ दो: यही रेखा \omterm{def:g6:symmetry:bisector}{लंब समद्विभाजक}
        है (हर कटान-बिंदु $A$ और $B$ से समान दूरी पर है)।
\end{enumerate}
\end{method}

\section{अभ्यास}

\begin{exercise}[$\star$]\label{exo:g6:symmetry:1}
जाल वाले काग़ज़ पर उतारो: एक खड़ी अक्ष $d$, उससे दो खाने बाईं ओर
एक बिंदु $A$, और पाँच खाने बाईं ओर तथा एक खाना ऊपर एक बिंदु $B$।
$d$ के पार उनके \omterm{def:g6:symmetry:def}{सममित} बिंदु $A'$ और $B'$ बनाओ।
\end{exercise}

\begin{exercise}[$\star$]\label{exo:g6:symmetry:2}
एक रेखा $d$ और उससे $3$ cm दूर एक बिंदु $M$ बनाओ (जाल वाले काग़ज़
पर नहीं)। \omterm{def:g6:shapes:triangles}{समकोण} पट्टिका और मापपट्टी से \cref{met:g6:symmetry:construct}
के अनुसार $d$ के पार $M$ का \omterm{def:g6:symmetry:def}{सममित} बिंदु $M'$ बनाओ। $M$ और
$M'$ कितनी दूर हैं?
\end{exercise}

\begin{exercise}[$\star$]\label{exo:g6:symmetry:3}
वर्णमाला के कौन-से बड़े अक्षर, सबसे सादे रूप में लिखने पर, खड़ी
\omterm{def:g6:symmetry:axis}{सममिति अक्ष} रखते हैं (जैसे A)? आड़ी अक्ष (जैसे B)?
दोनों?
\end{exercise}

\begin{exercise}[$\star$]\label{exo:g6:symmetry:4}
कितनी \omterm{def:g6:symmetry:axis}{सममिति अक्ष} हैं: \omterm{def:g6:shapes:triangles}{समबाहु त्रिभुज} की? (वर्ग न होने वाले)
आयत की? (वर्ग न होने वाले) \omterm{def:g6:shapes:quadrilaterals}{समचतुर्भुज} की? वर्ग की? वृत्त की?
हर आकृति अपनी अक्षों के साथ बनाओ।
\end{exercise}

\begin{exercise}[$\star$]\label{exo:g6:symmetry:5}
सही या ग़लत? ``आयत के विकर्ण उसकी \omterm{def:g6:symmetry:axis}{सममिति अक्ष} होते हैं।''
मोड़ने की दलील या चित्र से समझाओ।
\end{exercise}

\begin{exercise}[$\star$]\label{exo:g6:symmetry:6}
त्रिभुज $ABC$ को रेखा $d$ पर परावर्तित करके $A'B'C'$ बनाया गया।
$AB = 4$ cm, $\widehat{BAC} = 70^\circ$ और $ABC$ का \omterm{def:g4:measure:perimeter}{परिमाप}
$12$ cm दिया है: बिना कोई रचना किए $A'B'$,
$\widehat{B'A'C'}$ और $A'B'C'$ का \omterm{def:g4:measure:perimeter}{परिमाप} बताओ।
\end{exercise}

\begin{exercise}[$\star$]\label{exo:g6:symmetry:7}
$7$ cm का \omterm{def:g6:lines:objects}{रेखाखंड} $[AB]$ बनाओ और परकार से उसका \omterm{def:g6:symmetry:bisector}{लंब समद्विभाजक}
बनाओ (\cref{met:g6:symmetry:bisector})। उस पर कोई भी बिंदु $M$
चुनो और मापपट्टी से जाँचो कि $MA = MB$।
\end{exercise}

\begin{exercise}[$\star\star$]\label{exo:g6:symmetry:8}
किसी (समतल) बग़ीचे के बिंदुओं $A$ और $B$ पर दो पेड़ खड़े हैं।
फ़व्वारा कहाँ लगाया जाए कि वह दोनों पेड़ों से एक ही दूरी पर हो?
\emph{सारी} संभव जगहों का वर्णन करो।
\end{exercise}

\begin{exercise}[$\star\star$]\label{exo:g6:symmetry:9}
जाल वाले काग़ज़ पर अक्ष $d$ (जाल का $45^\circ$ वाला विकर्ण) और
उसकी एक ओर L-आकार की एक छोटी आकृति बनाओ। $d$ के पार उसकी
\omterm{def:g6:symmetry:def}{सममित} आकृति बनाओ। (हर \omterm{def:g5:solids:def}{शीर्ष} परावर्तित करो: $45^\circ$ वाली अक्ष
के लिए, अक्ष से $2$ खाने दाईं ओर का बिंदु उससे $2$ खाने ऊपर आ
जाता है।)
\end{exercise}

\begin{exercise}[$\star\star$]\label{exo:g6:symmetry:10}
केंद्र $O$ का एक वृत्त, $O$ से होकर एक रेखा $d$, और वृत्त पर एक
बिंदु $M$ बनाओ। $d$ के पार $M$ का \omterm{def:g6:symmetry:def}{सममित} बिंदु कहाँ है? समझाओ
कि परावर्तन वृत्त को उसी पर क्यों भेज देता है।
\end{exercise}

\begin{exercise}[$\star\star$]\label{exo:g6:symmetry:11}
समदूरी के गुण (\cref{prop:g6:symmetry:bisector}) से समझाओ कि
\cref{met:g6:symmetry:bisector} में चापों के दोनों कटान-बिंदु
सचमुच $[AB]$ के \omterm{def:g6:symmetry:bisector}{लंब समद्विभाजक} पर ही क्यों पड़ते हैं।
\end{exercise}

\begin{exercise}[$\star\star\star$]\label{exo:g6:symmetry:12}
बिंदु $A$ पर रखी बिलियर्ड की गेंद को सीधी दीवार $d$ से टकराकर
बिंदु $B$ तक पहुँचना है (दोनों दीवार की एक ही ओर हैं)। $B$ को
$d$ के पार परावर्तित करके $B'$ बनाओ, और \omterm{def:g6:lines:objects}{रेखाखंड} $[AB']$
खींचो। समझाओ कि टकराने की सबसे अच्छी जगह वहीं है जहाँ $[AB']$
दीवार को काटता है। (विचार: रास्ता $A \to P \to B$ और
$A \to P \to B'$ बराबर लंबे हैं।)
\end{exercise}

\section{समस्या: फ़व्वारा, टूटी थाली और चलता हुआ
दर्पण}

\begin{problem}[{सप्ताहांत समस्या --- त्रिभुज के तीनों लंब
समद्विभाजक एक ही बिंदु पर मिलते हैं: तीन बिंदुओं से होकर जाता
वृत्त, और दो दर्पण मिलकर क्या करते हैं}]
\label{pb:g6:symmetry:1}
\cref{exo:g6:symmetry:8} में फ़व्वारे को \emph{दो} पेड़ों से एक ही
दूरी पर होना था, और उत्तर जगहों की एक पूरी रेखा थी। यह समस्या
\emph{तीसरा} पेड़ लगाती है --- और पूरी रेखा सिमटकर एक ही, पूरी
तरह तय बिंदु रह जाती है। उस बिंदु में ज्यामिति का ख़ज़ाना छिपा
है: उससे तीन दिए हुए बिंदुओं से होकर जाने वाला इकलौता वृत्त बनता
है, टूटे ठीकरे से पूरी थाली दोबारा खड़ी होती है, और वह हर त्रिभुज
के पास होता है। आख़िरी भाग में दो दर्पणों का एक प्रयोग है जिसका
नतीजा अगले साल की ज्यामिति की घोषणा कर देता है।

\medskip
\noindent\textbf{भाग I --- दो पेड़, फिर इलाक़ों का
नक़्शा।}
\begin{enumerate}
  \item $AB = 6$ cm वाले दो बिंदु $A$ और $B$ बनाओ और परकार से
        $[AB]$ का \omterm{def:g6:symmetry:bisector}{लंब समद्विभाजक} बनाओ
        (\cref{met:g6:symmetry:bisector})। उस पर एक बिंदु
        $M$ चुनो और मापपट्टी से जाँचो कि $MA = MB$।
  \item \Cref{prop:g6:symmetry:bisector} एक साथ दो बातें कहती
        है: समद्विभाजक \emph{पर} पड़ा हर बिंदु $A$ और $B$ से
        समान दूरी पर है, और समान दूरी वाला हर बिंदु
        समद्विभाजक \emph{पर} है। इनमें से कौन-सी बात
        \cref{exo:g6:symmetry:8} के फ़व्वारे वाले सवाल का उत्तर
        देती है? समद्विभाजक पर \emph{न} पड़े बिंदु के बारे में
        प्रतिज्ञप्ति क्या कहती है?
  \item समद्विभाजक की उसी ओर एक बिंदु $M$ लो जिस ओर $A$ है।
        \omterm{def:g6:lines:objects}{रेखाखंड} $[MB]$ समद्विभाजक को किसी बिंदु $P$ पर काटता
        है। रास्ते $M \to P \to B$ की तुलना रास्ते
        $M \to P \to A$ से करो, और समझाओ कि $M$, $B$ की तुलना
        में $A$ के ज़्यादा पास क्यों है: समद्विभाजक शीट को एक
        ``$A$-ओर'' और एक ``$B$-ओर'' में बाँट देता है।
  \item किसी क़स्बे में दो विद्यालय $A$ और $B$ हैं। हर बच्चा पास
        वाले विद्यालय तक पैदल जाता है। $A$ और $B$ वाला एक छोटा
        नक़्शा बनाओ, और क़स्बे का वह ठीक-ठीक इलाक़ा रँगो जिसे
        विद्यालय $A$ सँभालता है। सीमा कौन-सी रेखा बनाती है?
  \item तीसरा विद्यालय $C$ खुलता है (उसे ऐसे रखो कि तीनों
        विद्यालय एक त्रिभुज बनाएँ)। $[AB]$ का \omterm{def:g6:symmetry:bisector}{लंब समद्विभाजक} और
        $[BC]$ का \omterm{def:g6:symmetry:bisector}{लंब समद्विभाजक} बनाओ, और उनके कटान-बिंदु को
        $O$ कहो। $O$ दूरियों की कौन-सी दो बराबरियाँ पूरी करता
        है, और क्यों?
\end{enumerate}

\medskip
\noindent\textbf{भाग II --- तीन पेड़, एक बिंदु, एक
वृत्त।}
\begin{enumerate}[resume]
  \item प्रश्न 5 से निकालो कि $OA = OC$ --- और इसलिए यह भी कि
        $O$, $[AC]$ के उस \omterm{def:g6:symmetry:bisector}{लंब समद्विभाजक} पर भी है जिसे तुमने
        कभी खींचा ही नहीं। निष्कर्ष निकालो: \emph{किसी त्रिभुज
        की भुजाओं के तीनों \omterm{def:g6:symmetry:bisector}{लंब समद्विभाजक} एक ही बिंदु से होकर
        जाते हैं।}
  \item समझाओ कि केंद्र $O$ और \omterm{def:g3:shapes:shapes}{त्रिज्या} $OA$ वाला वृत्त $B$ और
        $C$ से होकर भी क्यों जाता है: वही तीनों बिंदुओं से होकर
        जाने वाला \emph{इकलौता} वृत्त है --- त्रिभुज $ABC$ का
        \emph{परिवृत्त}\index{परिवृत्त}।
  \item $6$ cm, $5$ cm और $4$ cm भुजाओं वाला त्रिभुज बनाओ
        (\cref{met:g6:shapes:construct}), फिर उसका
        \omterm{pb:g6:symmetry:1}{परिवृत्त}। केंद्र पाने के लिए केवल \emph{दो} \omterm{def:g6:symmetry:bisector}{लंब
        समद्विभाजक} खींचना ही काफ़ी क्यों है?
  \item टूटी थाली: एक पुरातत्वविद् को ऐसा ठीकरा मिलता है जिसका
        केवल थाली की गोल कोर वाला हिस्सा साबुत है। ऐसी विधि
        बताओ जो मूल थाली का केंद्र और \omterm{def:g3:shapes:shapes}{त्रिज्या} वापस पा ले ---
        और उसे आज़माओ: परकार से एक चाप खींचो, केंद्र ``भूल''
        जाओ, चाप पर तीन बिंदु अंकित करो, और उसे दोबारा खड़ा करो।
  \item क्या कभी तीन बिंदुओं से होकर \emph{कोई} वृत्त न जाता
        हो? मान लो $A$, $B$, $C$ संरेख हैं। $[AB]$ और $[BC]$ के
        \omterm{def:g6:symmetry:bisector}{लंब समद्विभाजकों} के बारे में क्या कहा जा सकता है
        (\cref{prop:g6:lines:facts})? निष्कर्ष निकालो।
\end{enumerate}

\medskip
\noindent\textbf{भाग III --- हर जगह अक्ष, और चलता हुआ
दर्पण।}
\begin{enumerate}[resume]
  \item \Cref{ex:g6:symmetry:shapes} ने \omterm{def:g6:shapes:triangles}{समबाहु त्रिभुज} की $3$
        और वर्ग की $4$ अक्ष गिनी थीं। सम \omterm{def:g4:geometry:polygon}{पंचभुज} और सम \omterm{def:g4:geometry:polygon}{षट्भुज} के
        लिए गिनती का अनुमान लगाओ, और हर हाल में बताओ कि अक्ष
        कहाँ से गुज़रती हैं (भुजाओं की \omterm{def:g3:numbers:evenodd}{विषम संख्या} बनाम सम: दोनों
        स्थितियाँ अलग हैं)।
  \item काग़ज़ की शीट एक बार मोड़ो और दोनों परतों में से होकर कोई
        आकार काटो; फिर खोलो। जो भी काटो, छेद की एक \omterm{def:g6:symmetry:axis}{सममिति अक्ष}
        \emph{हमेशा} क्यों होती है, और कौन-सी रेखा वह अक्ष है?
  \item शीट दो बार मोड़ो --- दूसरा मोड़ पहले पर \omterm{def:g6:lines:perp}{लंब} --- काटो, और
        खोलो। छेद की कम से कम कितनी \omterm{def:g6:symmetry:axis}{सममिति अक्ष} होती हैं, और
        तुम्हारी कटाई की कितनी प्रतियाँ दिखती हैं? (काग़ज़ के
        बर्फ़-फाहे इसी विचार को और आगे मोड़कर बनते हैं।)
  \item चलता हुआ दर्पण: जाल वाले काग़ज़ पर एक खड़ी अक्ष $d_1$,
        उसकी बाईं ओर एक छोटा झंडा $F$, और $d_1$ से $4$ खाने
        दाईं ओर दूसरी खड़ी अक्ष $d_2$ बनाओ। $F$ को $d_1$ पर
        परावर्तित करो (प्रतिबिंब $F'$), फिर $F'$ को $d_2$ पर
        (प्रतिबिंब $F''$)। $F$ और $F''$ की तुलना करो: दिशा-क्रम
        वही है या दर्पण जैसा? आकृति कितने खाने खिसकी? अक्षों के
        बीच की दूरी से तुलना करो।
  \item वही प्रयोग $d_2$ को $d_1$ पर \emph{\omterm{def:g6:lines:perp}{लंब}} रखकर करो: झंडे
        को खड़ी अक्ष पर परावर्तित करो, फिर प्रतिबिंब को आड़ी अक्ष
        पर। बताओ कि $F''$, $F$ के और दोनों अक्षों के कटान-बिंदु
        के मुक़ाबले कैसे बैठता है। (तुमने अभी-अभी अगले साल का
        रूपांतरण खोज लिया: \cref{ch:g7:central}।)
\end{enumerate}
\end{problem}
